The model must be well-defined to be used in NMPC, in such a way 
to obtain reliable predictions of system behavior in response to various control actions.
Neural networks are universal approximators, and they can be used both for designing 
the entire dynamic model of a system and approximating the disturbances and noises 
a system can be affected.
Our purpose is to approximate the entire kinematic model of an unicycle with a neural model,
due to its low complexity. 
In this section we will start describinf briefly the kinematic model of a unicycle, and then 
we will describe the neural architecture used to approximate it.

\subsection{Unicycle kinematic model}
The kinematic model of the unicyle is a simplified representation of the movement of the vehicle
in temrs of position and orientation. It assumes there will be no lateral deformations. 
The state space is represented by its position and orientational
\begin{equation}
    [x, y, \theta]
\end{equation} 
while the input actions are represented by longitudinal and angular velocities the vehicle can 
assume
\begin{equation}
    [v, \omega]
\end{equation}
So, the motion of the unicycle can be described through the following equations:
\begin{eqnarray}
    \dot{x}(t) &=& v(t) \cos(\theta(t)) \\
    \dot{y}(t) &=& v(t) \sin(\theta(t)) \\
    \dot{\theta}(t) &=& \omega(t) \\
\end{eqnarray}

